% BEGIN GENERATED ARTIFACT: def:negation-normal-form-propositional-logic
\begin{definitionbox}{Definition (Negation Normal Form)}
\begin{definition}[Negation Normal Form]
\label{def:negation-normal-form-propositional-logic}
A formula is in \emph{negation normal form} if it uses only propositional
variables, negated propositional variables, conjunction, and disjunction, and
every negation occurs immediately before a propositional variable.
\end{definition}
\end{definitionbox}

\begin{remark*}[Standard quantified statement]
For a formula \(\varphi\in\WFF\),
\[
\varphi\text{ is in negation normal form}
\Longleftrightarrow
\varphi\text{ is built from literals using only }\land\text{ and }\lor.
\]
\end{remark*}

\begin{remark*}[Definition predicate reading]
\[
\operatorname{NegationNormalForm}(\varphi)
\]
means that \(\varphi\) has no conditionals, no biconditionals, and no negation
outside the literal level.
\end{remark*}

\begin{remark*}[Interpretation]
Negation normal form is the first cleanup stage for propositional formulas.
It does not force a formula to be a disjunction of terms or a conjunction of
clauses; it only restricts the available connectives and pushes negation down
to propositional variables.
\end{remark*}

\begin{remark*}[Examples]
The formulas \(P\), \(\neg P\), \(P\land(\neg Q\lor R)\), and
\((P\lor Q)\land(\neg R\lor S)\) are in negation normal form.
\end{remark*}

\begin{remark*}[Non-Examples]
The formulas \(P\to Q\), \(P\leftrightarrow Q\), and \(\neg(P\land Q)\) are not
in negation normal form. The first two use connectives not allowed in negation
normal form, and the last has a negation applied to a compound formula.
\end{remark*}

\begin{dependencies}
\begin{itemize}
  \item \hyperref[def:literal-propositional-logic]{Literal}
  \item \hyperref[def:well-formed-formula]{Well-Formed Formula}
  \item \hyperref[def:logical-equivalence-propositional-logic]{Logical Equivalence}
\end{itemize}
\end{dependencies}
% END GENERATED ARTIFACT: def:negation-normal-form-propositional-logic

% BEGIN GENERATED ARTIFACT: def:nnf-conversion-procedure-propositional-logic
\begin{definitionbox}{Definition (NNF Conversion Procedure)}
\begin{definition}[NNF Conversion Procedure]
\label{def:nnf-conversion-procedure-propositional-logic}
The \emph{NNF conversion procedure} converts a propositional formula to an
equivalent formula in negation normal form by the following steps:
\begin{enumerate}
\item eliminate biconditionals using biconditional equivalence laws;
\item eliminate conditionals using conditional equivalence laws;
\item eliminate double negations;
\item push negations inward using De Morgan laws.
\end{enumerate}
\end{definition}
\end{definitionbox}

\begin{remark*}[Standard quantified statement]
Starting with \(\varphi\in\WFF\), repeatedly apply the equivalences
\[
\alpha\leftrightarrow\beta
\equiv
(\alpha\to\beta)\land(\beta\to\alpha),
\qquad
\alpha\to\beta\equiv\neg\alpha\lor\beta,
\]
\[
\neg\neg\alpha\equiv\alpha,
\qquad
\neg(\alpha\land\beta)\equiv\neg\alpha\lor\neg\beta,
\qquad
\neg(\alpha\lor\beta)\equiv\neg\alpha\land\neg\beta,
\]
until negations occur only immediately before propositional variables and only
\(\land\) and \(\lor\) remain above the literal level.
\end{remark*}

\begin{remark*}[Definition predicate reading]
\[
\operatorname{NNFConversionProcedure}(\varphi,\psi)
\]
means that \(\psi\) is obtained from \(\varphi\) by eliminating
biconditionals and conditionals, removing double negations, and pushing
remaining negations inward.
\end{remark*}

\begin{remark*}[Interpretation]
The procedure is a syntax-directed rewrite process. Each rewrite step is backed
by a logical equivalence law, so the output formula has the same truth
conditions as the input formula.
\end{remark*}

\begin{dependencies}
\begin{itemize}
  \item \hyperref[thm:conditional-equivalence-laws-propositional-logic]{Conditional Equivalence Laws}
  \item \hyperref[thm:biconditional-equivalence-laws-propositional-logic]{Biconditional Equivalence Laws}
  \item \hyperref[thm:de-morgan-laws-propositional-logic]{De Morgan Laws for Propositional Logic}
  \item \hyperref[thm:basic-boolean-equivalence-laws-propositional-logic]{Basic Boolean Equivalence Laws}
\end{itemize}
\end{dependencies}
% END GENERATED ARTIFACT: def:nnf-conversion-procedure-propositional-logic

% BEGIN GENERATED ARTIFACT: thm:existence-negation-normal-form-propositional-logic
\begin{theorembox}{Theorem (Existence of Negation Normal Form)}
\begin{theorem}[Existence of Negation Normal Form]
\label{thm:existence-negation-normal-form-propositional-logic}
For every formula \(\varphi\in\WFF\), there exists a formula \(\psi\in\WFF\)
such that \(\psi\) is in negation normal form and
\[
\varphi\equiv\psi.
\]

\noindent\hyperref[prf:existence-negation-normal-form-propositional-logic]{Go to proof.}
\end{theorem}
\end{theorembox}

\begin{remark*}[Standard quantified statement]
\[
(\forall\varphi\in\WFF)(\exists\psi\in\WFF)
\bigl(\psi\text{ is in negation normal form}\land \varphi\equiv\psi\bigr).
\]
\end{remark*}

\begin{remark*}[Predicate reading]
\[
\operatorname{NegationNormalFormExists}(\varphi)
\]
means that \(\varphi\) has a logically equivalent representative in negation
normal form.
\end{remark*}

\begin{remark*}[Interpretation]
Every propositional formula can be rewritten so that negation appears only at
the variable level. This makes negation normal form a semantic normalizing step:
it changes shape while preserving logical equivalence.
\end{remark*}

\begin{dependencies}
\begin{itemize}
  \item \hyperref[def:negation-normal-form-propositional-logic]{Negation Normal Form}
  \item \hyperref[def:nnf-conversion-procedure-propositional-logic]{NNF Conversion Procedure}
  \item \hyperref[thm:replacement-of-logical-equivalents-propositional-logic]{Replacement of Logical Equivalents}
  \item \hyperref[thm:structural-induction-propositional-formulas]{Structural Induction for Propositional Formulas}
\end{itemize}
\end{dependencies}
% END GENERATED ARTIFACT: thm:existence-negation-normal-form-propositional-logic
